\chapter{Среда Yosis: лаборатория в браузере}
\label{ch:yosis}

В первой части мы рисовали схемы на бумаге и проверяли их в уме. Для ПЛИС
такой путь слишком долог: нужно описать схему текстом, проверить её, увидеть,
во что она превратилась, и только потом загрузить в кристалл. Все эти шаги мы
будем выполнять в одной программе --- \term{Yosis}. Она открывается в
браузере, ничего не требует устанавливать и показывает каждый этап пути,
который мы разобрали в главе~\ref{ch:flow}.

Внутри Yosis работают открытые инструменты: Yosys превращает текст в схему,
nextpnr размещает её на кристалле, Project Apicula упаковывает битстрим, а
Icarus Verilog проверяет поведение. Среда лишь собирает их вместе и рисует то,
что они вычислили.

\section{Карта среды}

Слева в Yosis находится колонка вкладок. Они идут в том же порядке, в каком
мы проходим путь от идеи до платы.

\begin{table}[H]
\centering
\caption{Вкладки Yosis и главы книги, которые к ним относятся}
\small
\begin{tabularx}{\textwidth}{l X l}
\toprule
\textbf{Вкладка} & \textbf{Что в ней делают} & \textbf{Глава} \\
\midrule
Обучение & Пошаговые уроки с вопросами и проверкой решения & --- \\
Код & Пишут описание схемы на Verilog & \ref{ch:verilog} \\
Выводы & Назначают порты схемы на ножки корпуса & \ref{ch:board} \\
RTL & Смотрят схему из логических элементов & \ref{ch:bridge} \\
Диаграммы & Подают сигналы и смотрят временные диаграммы & \ref{ch:sim} \\
Лаборатория & Подключают к плате детали и запускают схему «как на столе» & \ref{ch:examples} \\
Gowin & Видят LUT, триггеры и блоки сложения после синтеза & \ref{ch:inside} \\
Кристалл & Размещение и трассировка, отчёт о частоте, битстрим & \ref{ch:flow} \\
Подключение & Загрузка в плату по USB, UART-терминал & \ref{ch:flow} \\
\bottomrule
\end{tabularx}
\end{table}

\section{Новый проект}

Работа начинается с кнопки «Создать проект» в верхней строке. Среда
предлагает пять способов начать: написать Verilog с нуля, собрать схему
мышью из логических элементов, взять готовый пример, импортировать проект из
файла или открыть папку Gowin. Мы выберем «С нуля». Нужно только дать проекту
имя и указать имя верхнего модуля: обычно его называют \code{top}.

\shot{app-wizard}{Окно создания проекта: пять способов начать}{fig:y-wizard}

\section{Код}

Вкладка «Код» --- это редактор с подсветкой Verilog. Ошибки подчёркиваются
по мере набора, а сохранение происходит само. Для разбора возьмём проект
«бегущий огонь»: единица в шестиразрядном кольцевом регистре перемещается на
следующий светодиод восемь раз в секунду. Это схема из главы~\ref{ch:examples},
но порты названы так, как их называет плата: \code{clk}, \code{reset} и
\code{led\_n}.

\shot{app-code}{Редактор кода. Кнопка «Синтезировать» в углу запускает весь маршрут}{fig:y-code}

\begin{infobox}[Почему \code{led\_n}]
Суффикс \code{\_n} --- обычное соглашение: сигнал с активным \emph{низким}
уровнем. Светодиоды платы горят, когда на выводе 0, поэтому в схеме мы
подаём на них инвертированное значение кольца.
\end{infobox}

\section{RTL: схема, которую построил синтезатор}

Нажмём «Синтезировать». Меньше чем за секунду Yosys читает текст и строит
схему. Вкладка RTL (register-transfer level, «уровень регистровых передач»:
описание схемы как набора регистров и логики между ними) показывает её так,
как рисуют на бумаге: прямоугольники для регистров, условные знаки для
вентилей, кружки с плюсом для сумматоров. Это ровно то, что мы изучали в
первой части, только нарисованное автоматически.

\shot{br-adder-rtl}{RTL-схема на примере полного сумматора (глава~
ef{ch:bridge}): вентили XOR, AND и OR, подписи с номерами строк кода и справа инспектор}{fig:y-rtl}

Щелчок по элементу открывает справа инспектор: входы, выходы, параметры и
строку исходного текста, из которой элемент получился. Двойной щелчок
переносит в редактор. Колесо мыши меняет масштаб, перетаскивание двигает схему.

\begin{ideabox}
Среда показывает схему \emph{до} привязки к конкретному кристаллу. Это
«идеальная» схема из вентилей и триггеров. Внизу панели написано: «До
отображения на FPGA». Что из неё получится на настоящем кристалле, видно на
вкладке Gowin.
\end{ideabox}

\section{Выводы: от порта к ножке}

Порт \code{clk} в Verilog --- просто имя. Чтобы он заработал, нужно сказать,
к какой ножке корпуса подключён генератор 27~МГц. Для этого служит вкладка
«Выводы». Слева таблица портов вашей схемы, справа корпус GW2AR-18 QN88 с
подписями всех ножек. Выберите порт, щёлкните по ножке, нажмите «Назначить».
Среда запишет строку \code{IO\_LOC} в файл \code{board.cst} --- тот самый файл
ограничений, который мы разбирали в разделе~\ref{sec:cst}.

\shot{app-pins}{Планировщик выводов. Для портов с подходящими именами ножки подставлены сами}{fig:y-pins}

\section{Gowin: во что превратилась схема}

Вкладка Gowin показывает результат следующего шага --- отображения на
примитивы кристалла. Вместо абстрактных вентилей здесь LUT, триггеры DFF,
блоки арифметики ALU и буферы ввода-вывода. Если щёлкнуть по LUT, инспектор
выведет её таблицу истинности. Именно так наглядно видно утверждение из
первой части: логический элемент в ПЛИС --- это таблица в памяти.

\shot{br-adder-lut}{Тот же сумматор после отображения на примитивы Gowin: вместо вентилей остались две LUT, справа в инспекторе видна таблица истинности}{fig:y-gowin}

\section{Диаграммы: поведение во времени}

Схема нарисована, но работает ли она? Вкладка «Диаграммы» запускает
симулятор Icarus Verilog и рисует временные диаграммы сигналов
(рис.~\ref{fig:y-wave}). Для комбинационных схем есть кнопка «Перебрать
входы»: Yosis сам подаст все комбинации входных сигналов по очереди. Для
последовательных схем среда сама вырабатывает такт и короткий импульс сброса.
Подробнее о симуляции --- в главе~\ref{ch:sim}.

\shot{br-adder-wave}{Все восемь комбинаций входов полного сумматора: так проверяют таблицу истинности}{fig:y-wave}

\subsection{Сценарий: события, циклы и выбор сигналов}

Для схем с памятью одного перебора мало: нужно подать входы в нужные
моменты. Кнопка «Настроить сценарий» открывает панель с четырьмя вкладками.
\emph{Запуск} задаёт тактовый вход, сигнал сброса, период и число циклов.
\emph{Входы} хранит постоянные значения. На вкладке \emph{События} можно
добавлять шаги: «в периоде 5 вход \code{c5} равен 1, в периоде 6 --- снова 0».
Вкладка \emph{Проверки} позволяет записать ожидаемые значения, и Yosis
сравнит их с результатом симуляции.

Если сигналов много, меню «Сигналы» над диаграммой позволяет скрыть
служебные цепи и оставить только нужные, а кнопка «Показать весь интервал»
вписывает всю симуляцию в окно. Так построены диаграммы автомата и очереди
в главе~\ref{ch:fsmdesign}: все входные импульсы заданы на вкладке «События».

\section{Кристалл: размещение и битстрим}

Кнопка «Разместить и собрать» запускает nextpnr. Он решает, на какую именно
ячейку кристалла посадить каждую LUT и каждый триггер, и проводит между ними
соединения. На карте цветными квадратами видно, какая часть кристалла занята.
Рядом --- отчёт о частоте: для нашего примера он показывает запас более чем
в десять раз относительно 27~МГц.

\shot{app-chip}{Карта размещения на кристалле и результат nextpnr}{fig:y-chip}

После успешной сборки доступны файлы \code{.fs} (битстрим) и журнал.

\section{Лаборатория}

Прежде чем грузить схему в настоящую плату, её можно запустить на
виртуальном столе. Лаборатория показывает плату Tang Nano 20K с шестью
светодиодами и двумя кнопками. К гребёнке контактов можно подключить детали:
светодиоды, индикаторы, датчики, экраны, память. Схема исполняется в браузере
непрерывно, а детали ведут себя как настоящие, включая последовательные протоколы I$^2$C, SPI
и UART (глава~
ef{ch:uart}).

\shot{app-lab}{Бегущий огонь на виртуальной плате. Слева список из 68 деталей}{fig:y-lab}

\section{Подключение и обучение}

Вкладка «Подключение» загружает битстрим в плату прямо из браузера по
WebUSB: выбираете плату, выбираете файл, нажимаете «Загрузить в SRAM». Здесь
же UART-терминал: окно, в которое приходят символы, отправленные схемой
по последовательной линии, и из которого можно отправить символы в схему.

\shot{app-conn}{Загрузка битстрима в плату}{fig:y-conn}

А вкладка «Обучение» --- самостоятельный курс из тридцати одного урока для
тех, кто видит ПЛИС впервые: от «что такое бит» до собственных схем с
проверкой решений симулятором.

\shot{app-learn}{Курс «Твоя мастерская»: уроки с проверкой}{fig:y-learn}

\begin{taskbox}[Задание]
Создайте в Yosis проект с нуля, назовите модуль \code{top} и введите
описание с рис.~\ref{fig:y-code}. Нажмите «Синтезировать» и последовательно
пройдите вкладки RTL, Gowin, Кристалл, Лаборатория. На каждой найдите:
сколько триггеров в схеме (29: шесть в кольце, двадцать два в делителе и один для сигнала tick),
какую максимальную частоту обещает nextpnr и в каком порядке загораются
светодиоды.
\end{taskbox}
